\chapter{Reflection and Project Management}

\section{Summary of Progress}

During Semester One, the project has progressed from initial cleanroom
training to the definition of a clear and feasible experimental direction.
A full set of baseline depletion mode devices were fabricated early in the
project, providing functional devices and a reference dataset for later
comparison. In parallel, background reading was undertaken to establish
the physical principles underpinning enhancement-mode operation and to
identify suitable gate recess strategies compatible with the available
heterostructure and cleanroom processes. Together, this work has ensured
that Semester Two can focus on the fabrication of these devices and their
characterisation rather than exploratory research and preparation.

\section{Challenges and Adaptations}

The primary challenges encountered during this phase were practical
constraints on cleanroom availability and the inherent uncertainty
associated with selecting an appropriate enhancement-mode fabrication
route. These were mitigated through early completion of training,
close consultation with cleanroom staff, and a literature-led approach
to process selection. Periods of reduced experimental activity were
used productively for data analysis, tool development, and report
structuring, ensuring continued progress despite limited laboratory
access.

\section{Skills Developed}

The project has contributed to the development of a broad range of
technical and professional skills. These include hands-on experience
with III--V semiconductor fabrication, photolithography, wet etching,
metallisation, annealing, and wire bonding. In addition, significant
progress was made in experimental data handling through the development
of Python-based analysis tools, alongside improved proficiency in
technical writing, literature evaluation, and project planning using
structured timelines.

\section{Lessons for Semester Two}

The main lesson carried forward into Semester Two is the importance of
tight process control and realistic time planning when working within a
shared cleanroom environment. Experimental work will therefore be
structured around clearly defined weekly objectives, with sufficient
contingency for process iteration. Emphasis will be placed on systematic
variation, repeatability, and timely characterisation to ensure that
meaningful conclusions can be drawn within the remaining project
duration.
